% ============================================================
% preamble.tex — 기관 관리자 매뉴얼 공통 프리앰블
% xelatex + kotex + fontspec + tcolorbox
% ============================================================

\usepackage[a4paper,top=22mm,bottom=22mm,left=22mm,right=22mm,headheight=14pt,headsep=8mm,footskip=12mm]{geometry}
\usepackage{fontspec}
\usepackage{kotex}
\usepackage{xcolor}
\usepackage{graphicx}
\usepackage[hidelinks]{hyperref}
\usepackage{enumitem}
\usepackage{fancyhdr}
\usepackage{titlesec}
\usepackage{booktabs}
\usepackage{amssymb}
\usepackage{tcolorbox}
\tcbuselibrary{skins,breakable}

% ---------- 폰트 ----------
\setmainfont{NotoSerifKR-VF.ttf}[
  Path = C:/WINDOWS/Fonts/,
  UprightFeatures = {RawFeature={+axis={wght=500}}},
  BoldFeatures    = {RawFeature={+axis={wght=800}}}
]
\setsansfont{NotoSansKR-VF.ttf}[
  Path = C:/WINDOWS/Fonts/,
  UprightFeatures = {RawFeature={+axis={wght=500}}},
  BoldFeatures    = {RawFeature={+axis={wght=800}}}
]
\setmonofont{Consolas}
% kotex 의 \setmainhangulfont / \setsanshangulfont 로 한글 처리
\setmainhangulfont{NotoSerifKR-VF.ttf}[
  Path = C:/WINDOWS/Fonts/,
  UprightFeatures = {RawFeature={+axis={wght=500}}},
  BoldFeatures    = {RawFeature={+axis={wght=800}}}
]
\setsanshangulfont{NotoSansKR-VF.ttf}[
  Path = C:/WINDOWS/Fonts/,
  UprightFeatures = {RawFeature={+axis={wght=500}}},
  BoldFeatures    = {RawFeature={+axis={wght=800}}}
]

% ---------- 색상 ----------
\definecolor{kfgreen}{HTML}{2f7a3e}
\definecolor{kfgreensoft}{HTML}{e8f4ec}
\definecolor{kfblue}{HTML}{1e6fdc}
\definecolor{kfbluesoft}{HTML}{e7f0fc}
\definecolor{kforange}{HTML}{e07a1a}
\definecolor{kforangesoft}{HTML}{fdf2e3}
\definecolor{kfgray}{HTML}{4a5568}
\definecolor{kfgrayline}{HTML}{cbd5e0}

% ---------- 챕터/섹션 타이틀 ----------
\titleformat{\chapter}[display]
  {\sffamily\Huge\bfseries\color{kfgreen}}
  {\fontsize{14}{16}\selectfont CHAPTER \thechapter}
  {6pt}
  {\Huge}
\titlespacing*{\chapter}{0pt}{-20pt}{20pt}

\titleformat{\part}[display]
  {\sffamily\Huge\bfseries\centering\color{kfgreen}}
  {\Large PART \thepart}
  {12pt}
  {\Huge}

\titleformat{\section}{\sffamily\Large\bfseries\color{kfgreen}}{}{0pt}{}
\titleformat{\subsection}{\sffamily\large\bfseries\color{kfgray}}{}{0pt}{}

% ---------- 헤더/푸터 ----------
\pagestyle{fancy}
\fancyhf{}
\fancyhead[L]{\sffamily\small\color{kfgray}국어농장v2 기관 관리자 매뉴얼}
\fancyhead[R]{\sffamily\small\color{kfgray}\leftmark}
\fancyfoot[C]{\sffamily\small\thepage}
\renewcommand{\headrulewidth}{0.4pt}
\renewcommand{\footrulewidth}{0pt}
\renewcommand{\chaptermark}[1]{\markboth{#1}{}}

% ---------- 박스 매크로 ----------
% 1) summarybox — 한눈에 요약 (초록)
\newtcolorbox{summarybox}[1]{
  enhanced,
  colback=kfgreensoft,
  colframe=kfgreen,
  coltitle=white,
  colbacktitle=kfgreen,
  title={\sffamily\bfseries #1},
  fonttitle=\bfseries,
  boxrule=0.6pt,
  arc=4pt,
  left=10pt, right=10pt, top=8pt, bottom=8pt,
  breakable
}

% 2) screenbox — 화면 구조 다이어그램 컨테이너 (회색)
%    안에는 booktabs 표 또는 tcolorbox 카드 사용 (verbatim ASCII 금지 — 한글 폭 정렬 깨짐)
\newtcolorbox{screenbox}[1]{
  enhanced,
  colback=white,
  colframe=kfgrayline,
  coltitle=kfgray,
  colbacktitle=white,
  title={\sffamily\bfseries\color{kfgray} #1},
  boxrule=0.5pt,
  arc=2pt,
  left=8pt, right=8pt, top=6pt, bottom=6pt,
  breakable
}

% 화면 카드 — 사이드바·메인 영역 등을 작은 박스로 표현
\newtcolorbox{screencard}[1]{
  enhanced,
  colback=kfbluesoft!50,
  colframe=kfblue,
  coltitle=kfblue,
  fonttitle=\sffamily\bfseries\small,
  title=#1,
  boxrule=0.5pt,
  arc=2pt,
  left=6pt, right=6pt, top=4pt, bottom=4pt,
  fontupper=\sffamily\small,
  width=\linewidth,
  before skip=2pt,
  after skip=2pt
}

% 탭 바 — 가로로 늘어선 탭 표현
\newcommand{\screentabs}[1]{%
  \par\noindent
  \tcbox[colback=kfgreensoft,colframe=kfgreen,boxrule=0.5pt,arc=2pt,left=4pt,right=4pt,top=2pt,bottom=2pt,nobeforeafter,box align=center]{%
    \sffamily\bfseries\small #1%
  }\par\smallskip%
}

% 활성 탭 강조
\newcommand{\activetab}[1]{\colorbox{kfgreen}{\textcolor{white}{\sffamily\bfseries\small\,#1\,}}}
\newcommand{\inacttab}[1]{\fcolorbox{kfgrayline}{white}{\sffamily\small\,#1\,}}

% 3) tipbox — 알아두면 좋은 점 (주황)
\newtcolorbox{tipbox}[1]{
  enhanced,
  colback=kforangesoft,
  colframe=kforange,
  coltitle=white,
  colbacktitle=kforange,
  title={\sffamily\bfseries 💡 #1},
  boxrule=0.6pt,
  arc=4pt,
  left=10pt, right=10pt, top=8pt, bottom=8pt,
  breakable
}

% 4) warnbox — 주의 (빨강)
\newtcolorbox{warnbox}[1]{
  enhanced,
  colback=red!5!white,
  colframe=red!70!black,
  coltitle=white,
  colbacktitle=red!70!black,
  title={\sffamily\bfseries ⚠ #1},
  boxrule=0.6pt,
  arc=4pt,
  left=10pt, right=10pt, top=8pt, bottom=8pt,
  breakable
}

% ---------- 캡처 placeholder ----------
% 사용법: \kfShot{파일명}{캡션}
% images/파일명.png 가 있으면 그걸 표시, 없으면 placeholder 박스
\newcommand{\kfShot}[2]{%
  \begin{figure}[!ht]
    \centering
    \IfFileExists{images/#1.png}{%
      \includegraphics[width=0.95\linewidth,keepaspectratio]{images/#1.png}%
    }{%
      \fbox{\parbox[c][5cm][c]{0.9\linewidth}{\centering\itshape\color{kfgray}(캡처 자리: #2 — \texttt{images/#1.png} 자동 생성 예정)}}%
    }
    \caption{#2}
  \end{figure}
}

% ---------- 표지 ----------
\newcommand{\kfTitlePage}{
  \begin{titlepage}
    \centering
    \vspace*{2cm}
    {\sffamily\fontsize{18}{22}\selectfont\color{kfgray} 국어농장 v2}\\[1.2em]
    {\sffamily\fontsize{36}{44}\selectfont\bfseries\color{kfgreen} 기관 관리자 매뉴얼}\\[1em]
    {\sffamily\fontsize{14}{18}\selectfont\color{kfgray} ORG\_ADMIN Operations Guide}\\[2em]
    \rule{0.4\textwidth}{0.6pt}\\[1em]
    {\sffamily\large\color{kfgray} 학원·교습소·교실 운영자를 위한 14개 화면 통합 안내}\\[6em]
    {\sffamily\small\color{kfgray} 발행: 국어농장 \quad / \quad 빌드일: \today}
    \vfill
  \end{titlepage}
}
